\scanpage{38}
2) Premières propriétés [\unsure{concernant} $\pi_0,\pi_1,\mathrm{Pic}$]

\uline{Proposition} Supposons que $H^0(X,\underline{O}_X)\longrightarrow H^0(X_{/Y},O_{X/Y})$ soit injectif [bijectif] alors $\pi_0(Y)\longrightarrow\pi_0(X)$ est surjectif [bijectif]. En particulier $(LG)_{-1}$ [resp$^{\mathrm{t}}$ $(LG)_0$] implique que $\pi_0(Y)\longrightarrow\pi_0(X)$ est surjectif [bijectif].

\uline{Proposition} \unsure{1|,} Supposons que $(X,Y)$ satisfasse \deleted{\unsure{L}} $LG$ (resp. $LG$ stricte). \inserted{soit $a$ un pt géom de $Y$ \unsure{supposé connexe}}, Alors $\pi_1(Y)\longrightarrow\pi_1(X)$ est surjectif [resp.
\[
\pi_1(Y,\deleted{\unsure{a}})\longrightarrow
\varprojlim_{\substack{U\ \text{voisinage}\\\text{ouvert de }Y}}\pi_1(U,\deleted{\unsure{a}})
\]
est un hom \deleted{surjectif [resp.} bijectif] ].

Résulte \unsure{des|de} faits plus généraux

\uline{\deleted{Corollaire}\inserted{Lemme}} Supposons $LG$. \deleted{\unsure{4} Alors pour tt} \inserted{\unsure{Soit}} $F$ loc libre sur $X$, (i) tte structure d'alg. \deleted{[resp. associative, \unsure{à élt unité|resp. unitaire}, resp. comm.]} \deleted{\unsure{comm}} sur $\hat F$ provient d'une structure d'algèbre unique sur $F$, qui est ass (resp unitaire, resp. comm) sss celle de $\hat F$ l'est. \deleted{\deleted{Pour que} \inserted{(ii) les décomp$^{\mathrm{ns}}$ de} $F$ \deleted{soit} \inserted{en} \deleted{\unsure{un}} produits de \inserted{\unsure{deux}} algèbres non nulles correspondent aux décompositions analogues de $\hat F$ [\unsure{non|une|en} \unsure{triviale|biuniv.} sss \deleted{\unsure{X}} celle de $F$ l'est].} Donc si $X'=\operatorname{Spec}(F)$, alors $X'$ connexe sss $\hat X'_{/Y}$ connexe

% margin note written sideways
\margin{(iii) Un hom \deleted{\unsure{de faisceaux}} \inserted{$F\longrightarrow G$ de Modules munis de} structure d'alg. est un hom d'alg, sss $\unsure{\vec F|F|\hat F}_{/Y}\longrightarrow G_{/Y}$ l'est.}

% slanted
\margin{\deleted{\uline{Lemme} [\unsure{Linéariser|L'énoncer}, donner un lemme \unsure{préliminaire} valable chaque fois qu'on a un foncteur \uline{\unsure{pl. fidèle|plt. fidèle|v.l. fidèle}} de catégories additives]}}
% archive sheet number: 37
